\documentclass[12pt]{article}
\usepackage[utf8]{inputenc}
\usepackage[T1]{fontenc}
\usepackage[margin=1in]{geometry}
\usepackage{setspace}
\usepackage{booktabs}
\usepackage{longtable}
\usepackage{array}
\usepackage{hyperref}

\title{Methods and Discussion Sections}
\author{Anonymous Author(s)}
\date{}

\begin{document}

\maketitle

\section{Methods}

\subsection{Research Design}
This study employed a qualitative longitudinal design to investigate how Azerbaijani-Iranian trilingual scholars use generative AI (GenAI) in academic writing and how such use reshapes authorship, agency, and epistemic responsibility. A qualitative approach was selected because the research sought to capture not only observable patterns of tool use, but also participants' reflexive accounts of decision-making, sense-making, and ownership during writing. The longitudinal orientation enabled the study to follow writing practices across multiple stages of composition, revision, and translation, thereby providing a process-sensitive account of human--AI collaboration.

The study was conceptually grounded in the notion of \emph{symbiotic agency}, defined here as a distributed but accountable form of authorship in which human writers and GenAI systems jointly contribute to the production of academic text while the human scholar retains epistemic responsibility for final decisions. The design was informed by scholarship on writing process theory, multilingual academic writing, and human--AI collaboration.

\subsection{Participants and Context}
The participants were eight Azerbaijani-Iranian scholars working across a range of disciplines and writing in Azerbaijani, Persian, and English. All participants had prior experience with academic writing in at least two languages, and all reported some degree of engagement with GenAI tools for tasks such as brainstorming, drafting, translation, paraphrasing, revision, and argument development. To protect anonymity, participants are identified using pseudonyms (P01--P08). Table~\ref{tab:participants} summarizes their disciplinary background, academic rank, and reported AI experience.

The study context is significant because Azerbaijani-Iranian scholars often navigate complex multilingual writing demands. For many of them, English functions as an additional academic language rather than a primary language of scholarly production. In such settings, GenAI is not merely a convenience tool; it can function as a linguistic mediator, cognitive scaffold, and epistemic partner.

\subsection{Data Collection}
Data were collected over an eight-month period using a combination of methods designed to capture both real-time interaction and retrospective reflection. First, participants completed writing-aloud sessions in which they verbalized their thoughts while drafting or revising academic text with the assistance of GenAI. These sessions provided fine-grained insight into the immediate reasoning behind prompt formulation, acceptance or rejection of AI outputs, and movement between languages.

Second, prompt interaction logs were collected to document the sequence, content, and functional purpose of prompts used during academic writing. These logs included prompts related to conceptual clarification, linguistic refinement, structural organization, and factual verification. In total, the dataset included more than 400 AI prompts across the eight participants.

Third, stimulated recall interviews were conducted to elicit participants' explanations of their writing choices after the fact. During these interviews, participants reviewed selected excerpts from their own AI interactions and reflected on why specific prompts were used, how AI suggestions were evaluated, and when human judgment overrode machine-generated output.

In addition, a reflexive AI-authorship agency scale was used as a supplementary instrument to capture participants' perceived control, dependence, and responsibility in relation to AI-supported writing. The scale was not intended as a psychometric diagnostic measure but as a structured reflection tool aligned with the study's qualitative aims.

\subsection{Data Analysis}
Data were analyzed using thematic analysis with an iterative coding process. The analysis began with repeated reading of transcripts, prompt logs, and reflective notes to identify recurring patterns in how participants framed GenAI use. Initial codes were generated for linguistic assistance, conceptual support, structural editing, cognitive offloading, epistemic monitoring, and algorithmic dependency. These codes were then compared across participants and across writing stages to identify higher-order themes.

Particular attention was given to moments in which participants either retained critical control over the text or ceded interpretive authority to the AI system. This allowed for the differentiation between productive symbiosis and problematic dependency. The final analytic framework resulted in four interrelated dimensions of symbiotic agency: linguistic empowerment, cognitive offloading, epistemic monitoring, and dependency risk.

To enhance trustworthiness, the analysis followed principles of credibility, transferability, dependability, and confirmability. Multiple data sources were triangulated, and analytic decisions were documented through memoing. The iterative comparison of writing-aloud data, prompt logs, and stimulated recall interviews strengthened the validity of thematic interpretations.

\subsection{Ethical Considerations}
Ethical approval was obtained in accordance with the norms of qualitative research involving human participants. All participants provided informed consent and were assured of confidentiality, anonymity, and the right to withdraw from the study at any time. Identifying information was removed from transcripts and logs, and pseudonyms were used in all reporting.

Because the study involved AI-mediated writing practices, special attention was paid to the ethics of authorship, transparency, and data privacy. Participants retained ownership of their writing, and no raw personal data were shared outside the research context. The study did not evaluate participants' work for correctness or productivity; rather, it examined the processes through which academic text is co-constructed in human--AI interaction.

\subsection{Methodological Limitations}
The study has several limitations. First, the sample size was small and context-specific, focusing on eight Azerbaijani-Iranian scholars. While appropriate for an in-depth qualitative inquiry, the findings should not be generalized statistically. Second, the rapid evolution of GenAI tools means that patterns observed during the study period may shift as interfaces and model capabilities change. Third, writing-aloud and stimulated recall methods may not capture all tacit aspects of writing behavior. Finally, because the study focused on multilingual academic writing in a specific sociolinguistic context, the findings may require adaptation before being transferred to other settings.

\subsection{Summary}
Overall, the methodological design enabled a detailed examination of how GenAI functions within multilingual academic writing as a tool of mediation, scaffolding, and potential dependency. By combining writing-aloud sessions, prompt logs, and stimulated recall interviews, the study captured both observable interaction patterns and participants' internal rationales, thereby providing a robust basis for theorizing symbiotic agency.

\section{Discussion}

\subsection{GenAI as a Mediator of Multilingual Academic Writing}
The findings demonstrate that GenAI is not simply an external writing aid but an active mediator in multilingual academic composition. For Azerbaijani-Iranian scholars, especially those writing in English as an additional language, AI tools supported lexical retrieval, grammatical correction, syntactic restructuring, and translation. In this sense, GenAI extended participants' linguistic repertoire and reduced the burden of producing academically acceptable text in a non-dominant language. This aligns with prior work showing that digital tools can redistribute the labor of writing across human and machine actors.

However, the study also shows that linguistic assistance alone does not capture the full significance of GenAI use. Participants did not merely ask the system to ``correct'' their writing; they used it to negotiate meaning, test formulations, and explore the boundaries of argumentation. Thus, GenAI functioned as a mediational means that shaped not only the surface form of writing but also the cognitive and rhetorical processes underlying composition.

\subsection{Symbiotic Agency and the Redistribution of Cognitive Labor}
One of the central contributions of this study is the concept of \emph{symbiotic agency}. The data suggest that authorship in AI-mediated writing is neither fully human nor fully automated. Instead, it is distributed across the scholar, the prompt, the interface, and the model's generated response. Yet this distribution does not eliminate human accountability. Rather, it reconfigures agency as a negotiated and recursive process in which the writer initiates, evaluates, revises, and authorizes the final text.

The notion of cognitive offloading was especially important in understanding this process. Participants frequently delegated lower-level tasks to GenAI, such as sentence polishing or paraphrase generation, in order to preserve cognitive resources for higher-order tasks like argument development and conceptual framing. When used strategically, this offloading appeared to support productivity and reduce linguistic friction. However, when overused or insufficiently monitored, it risked replacing critical engagement with surface fluency.

\subsection{Epistemic Monitoring as a Core Mechanism}
A key finding of the study is that epistemic monitoring is the mechanism that distinguishes productive human--AI collaboration from dependency-producing use. Participants who consistently verified AI-generated suggestions, cross-checked claims, and revised outputs according to disciplinary knowledge tended to maintain strong authorship control. By contrast, participants who accepted fluent output with minimal scrutiny were more vulnerable to the fluency--authority trap, in which polished language was mistaken for conceptual accuracy or intellectual ownership.

Epistemic monitoring therefore emerges as a crucial component of AI literacy in academic writing. It involves not only checking factual correctness but also evaluating whether the AI-supported text aligns with the writer's intended argument, disciplinary conventions, and ethical responsibilities. In this sense, epistemic monitoring is not a peripheral habit but a central aspect of scholarly authorship in the age of GenAI.

\subsection{Implications for Multilingual Scholars}
The study has important implications for multilingual academics, particularly those working in contexts where English is an additional language. GenAI can reduce inequalities in access to linguistic resources by offering immediate support for drafting and revision. This may be especially valuable for scholars who have strong conceptual expertise but face barriers in expressing that expertise in English.

At the same time, the findings caution against viewing AI as a neutral equalizer. The benefits of GenAI are unevenly distributed and depend on users' critical competence, prior writing experience, and access to quality prompting practices. Scholars with stronger epistemic monitoring skills appear better positioned to benefit from AI without becoming dependent on it. Thus, rather than replacing linguistic competence, GenAI shifts the focus toward evaluative and metacognitive competence.

\subsection{Theoretical Contribution}
The study contributes to writing studies, applied linguistics, and human--AI collaboration research in several ways. First, it extends process-oriented models of writing by showing how AI systems participate in drafting, revising, and translating academic text. Second, it contributes to multilingual scholarship by illustrating how writers working across Azerbaijani, Persian, and English strategically mobilize GenAI to manage translanguaging and linguistic mediation. Third, it advances theory on human--AI collaboration by proposing symbiotic agency as a framework that integrates empowerment, offloading, monitoring, and dependency risk.

Importantly, the study does not portray GenAI use as inherently beneficial or harmful. Instead, it emphasizes relational conditions: the quality of engagement, the presence of critical monitoring, and the writer's willingness to retain epistemic responsibility. This nuanced position moves beyond simplistic debates about whether AI ``helps'' or ``hurts'' writing.

\subsection{Pedagogical and Institutional Implications}
The findings suggest that higher education institutions should move beyond restrictive or purely detection-based approaches to AI use. Rather than treating GenAI as a threat to academic integrity, universities should develop policies and pedagogies that teach responsible, transparent, and reflective AI use. This includes training students and scholars to formulate effective prompts, evaluate outputs critically, document AI assistance where appropriate, and maintain authorship integrity.

Supervisors and writing instructors also have a role to play in helping multilingual scholars use AI ethically and productively. This may involve explicit discussion of when AI-supported language editing is acceptable, how to distinguish editing from ghostwriting, and how to preserve the writer's own voice. Institutional guidelines should therefore focus less on prohibition and more on responsible participation.

\subsection{Ethical Tensions}
The study highlights several ethical tensions associated with GenAI-mediated academic writing. One is linguistic homogenization: as AI systems tend to produce standardized academic prose, there is a risk that distinctive stylistic or cultural voice may be flattened. Another is unequal access: not all scholars have the same technological literacy or familiarity with prompting strategies, which may create new forms of inequality. A third is epistemic accountability: if AI-generated text is used without careful scrutiny, responsibility for claims becomes ambiguous.

These tensions reinforce the need for transparent authorship practices. Scholars should remain answerable for the intellectual content of their work even when the writing process is distributed across human and machine participants. In this sense, ethical AI use is not merely a procedural matter but part of scholarly integrity itself.

\subsection{Conclusion of the Discussion}
Overall, the discussion supports the argument that GenAI has entered academic writing as a mediating infrastructure rather than a mere convenience tool. For trilingual Azerbaijani-Iranian scholars, the technology opens new possibilities for expression, translation, and composition, but it also increases the importance of monitoring, judgment, and accountability. The proposed symbiotic agency model captures this dual reality by recognizing both the enabling and risky dimensions of human--AI writing collaboration.

Rather than asking whether AI should be used in academic writing, the more productive question is how writers can engage with AI in ways that preserve critical authorship while expanding linguistic and cognitive capacity. The answer, this study suggests, lies in a disciplined form of symbiosis: collaborative, but not passive; assisted, but not surrendered.

\end{document}
